\begin{helper}[Appending a certified alternating segment]
Let \(G\) be a finite loopless multigraph, let \(S\) be an ordered
tree-sequence rooted at \(e_0\), and let \(\mu\) be a partial
\(k\)-edge-coloring.  Suppose \(\alpha,\beta\in\overline\mu(S)\).
Let \(m\geq1\), and suppose we are given injective maps
\[
 v:\{0,\ldots,m\}\longrightarrow V,
 \qquad f:\{0,\ldots,m-1\}\longrightarrow E
\]
with the following properties.  The vertex \(v_0\) belongs to \(V(S)\),
every \(v_{i+1}\) lies outside \(V(S)\), every exact edge \(f_i\) is
distinct from every displayed edge of \(S\), and
\[
 \operatorname{ends}(f_i)=\{v_i,v_{i+1}\},\qquad
 \mu(f_i)\in\{\alpha,\beta\}.                              \tag{1}
\]
Fix a vertex \(y\), and also suppose that every vertex-node \(v_i\) and
edge-node \(f_i\) lies in the \((\alpha,\beta)\)-component through \(y\).

Then appending
\((f_0,v_1),\ldots,(f_{m-1},v_m)\) produces a tree-sequence \(S'\) such
that \(S\) is its exact initial segment,
\(|S'|=|S|+m>|S|\), and \(S'\) is reached from \(S\) by a
\noderef{TashkinovAugmentationSequence}.  Its last vertex is \(v_m\).
Every appended edge and vertex has the stated component membership and
color, every appended vertex lies outside \(V(S)\), and for
\(|S|\leq r<|S'|\),
\[
 \operatorname{ends}(S'.\operatorname{edge}(r))=
 \begin{cases}
  \{v_0,S'.\operatorname{vertex}(r+1)\},&r=|S|,\\
  \{S'.\operatorname{vertex}(r),S'.\operatorname{vertex}(r+1)\},&r>|S|.
 \end{cases}                                                \tag{2}
\]
\end{helper}
\begin{proof}
Define the enlarged vertex map on the ordinal sum of the old
\(|S|+1\) positions and the \(m\) new positions by retaining every old
entry and sending the new position numbered \(i\) to \(v_{i+1}\).
Define the enlarged edge map similarly, retaining the old \(|S|\) entries
and sending the new position numbered \(i\) to \(f_i\).  These maps have
domains of sizes \(|S|+m+1\) and \(|S|+m\), respectively.

The old part of the vertex map is injective by the
\noderef{TreeSequence} axioms for \(S\); the new part is injective because
\(v\) is injective; and the two images are disjoint because every
\(v_{i+1}\) is outside \(V(S)\), whose image description is
\noderef{TreeSequenceVertexSet}.  Hence the enlarged vertex map is
injective.  The same three-part argument applies to edges: the old map is
injective by \noderef{TreeSequence}, the new map is injective by hypothesis,
and old and new exact edge identities are disjoint by the freshness
hypothesis.  Thus the enlarged edge map is injective.

The first displayed edge remains the old first edge \(e_0\), because
\(|S|\geq1\).  Every old attachment equation is inherited literally from
\(S\).  For the first new edge, (1) attaches \(f_0\) to the earlier old
vertex \(v_0\in V(S)\); by \noderef{TreeSequenceVertexSet}, choose its
displayed index in \(S\).  For \(i>0\), equation (1) attaches \(f_i\) to
the immediately preceding new vertex \(v_i\).  All these attachment
indices precede the corresponding new vertex.  These facts construct all
fields of a \noderef{TreeSequence} \(S'\).  Its definitions give
\(|S'|=|S|+m\), make \(S\) an exact
\noderef{TreeSequenceInitialSegment}, identify its last vertex with
\(v_m\), and translate (1) entry by entry into (2).  Since \(m\geq1\),
the length inequality is strict.  The outside-tree, component-membership,
and color conclusions are immediate from the definitions of the appended
maps and the corresponding hypotheses.

It remains to construct the augmentation sequence.  For
\(0\leq t\leq m\), let \(S_t\) be the same construction using only the
first \(t\) new pairs.  Thus \(S_0=S\), \(S_m=S'\), and
\(|S_{t+1}|=|S_t|+1\).  The map definitions give an exact
\noderef{TreeSequenceInitialSegment} from \(S_t\) to \(S_{t+1}\).

Unfolding \noderef{TreeMissingColors}, choose a vertex of \(S\) missing
\(\alpha\) and a vertex of \(S\) missing \(\beta\).  Every \(S_t\)
retains all old displayed vertices literally, so the same two vertices,
with their unchanged missing-color facts under the fixed coloring \(\mu\),
show that both \(\alpha\) and \(\beta\) belong to
\(\overline\mu(S_t)\).  By (1), the new last edge in the step
\(S_t\to S_{t+1}\) has one of those two colors.  The one-edge length
increase, exact initial-segment equality, its stated color, and the retained
missing-color witness are precisely the clauses of
\noderef{TashkinovAugmentation}.  Hence every adjacent pair
\((S_t,S_{t+1})\) is one augmentation.  Finite induction on \(t\), using
the reflexive-transitive-closure definition in
\noderef{TashkinovAugmentationSequence}, concatenates these \(m\) steps
from \(S\) to \(S'\).  This proves the remaining assertion.
\end{proof}
